% !TEX program = xelatex
% =====================================================================
%  OVERLEAF: Menu -> Compiler -> XeLaTeX   (dong magic comment tren cung
%  thuong du, nhung dat trong Menu cho chac).
%
%  Cau truc thu muc:
%      bao_cao.tex
%      hinh/*.pdf        <- 23 file do bai_lam.ipynb tu sinh
%
%  Hinh nao chua co thi macro \hinh se ve mot o trong thay cho no,
%  file van bien dich duoc -> khong bi ket khi thieu hinh.
% =====================================================================
\documentclass[aspectratio=169, 11pt]{beamer}

\usetheme{default}
\usecolortheme{seahorse}
\setbeamertemplate{navigation symbols}{}
\setbeamertemplate{footline}[frame number]
\setbeamercolor{structure}{fg=blue!55!black}

\usepackage{fontspec}
\setsansfont{TeX Gyre Heros}    % co san tren Overleaf, day du dau tieng Viet
% \usepackage{polyglossia}\setmainlanguage{vietnamese}   % chi anh huong ngat tu, khong bat buoc

\usepackage{amsmath, amssymb}
\usepackage{booktabs}
\usepackage{graphicx}

% \hinh{ten.pdf} — thieu file thi ve o trong thay vi bao loi
\newcommand{\hinh}[1]{%
  \IfFileExists{hinh/#1}%
    {\includegraphics[width=\linewidth]{hinh/#1}}%
    {\fbox{\parbox[c][3.2cm][c]{0.95\linewidth}{\centering\footnotesize\ttfamily
      #1\\[2pt]\rmfamily (chưa có --- chạy notebook rồi bỏ vào thư mục \texttt{hinh/})}}}}

\title{Tối ưu hoá hàm log-loss\\bằng các thuật toán gradient}
\subtitle{Santander Product Recommendation}
\author{Nhóm ...}
\date{\today}

\begin{document}

\frame{\titlepage}

\begin{frame}{Nội dung}
  \tableofcontents
\end{frame}

% =====================================================================
\section{Bài toán và dữ liệu}

\begin{frame}{Bài toán}
  Hồi quy logistic dự đoán khách hàng \textbf{mua thêm} sản phẩm ở tháng $t$.
  \[
    F(w)=\underbrace{\frac1n\sum_{i=1}^n\Big[\log\big(1+e^{z_i}\big)-y_iz_i\Big]}_{\ell(w)\ \text{(log-loss)}}
    \;+\;\underbrace{\frac{\lambda}{2}\|w\|_2^2}_{\text{Ridge}}
    \quad\text{hoặc}\quad \ell(w)+\underbrace{\lambda_1\|w\|_1}_{\text{Lasso}},
    \qquad z=Xw
  \]
  \begin{itemize}
    \item Nhãn: $y_i=\mathbb{1}[\text{có ở tháng }t]\cdot\mathbb{1}[\text{chưa có ở tháng }t-1]$
    \item Gradient phần trơn: $\nabla\ell(w)=\frac1nX^\top\big(\sigma(Xw)-y\big)$
    \item Hằng số Lipschitz: $L=\dfrac{\sigma_{\max}(X)^2}{4n}+\lambda$ \quad(vì $\sigma'(z)\le\frac14$)
          $\;\Rightarrow\; L=6.25$, bước an toàn $1/L=0.160$
  \end{itemize}
  \vfill
  {\footnotesize Quy mô: train $n=466\,243$ dòng, validation $75\,839$ dòng, $d=161$ cột
  (160 đặc trưng + chặn). Tỉ lệ dương (sản phẩm \texttt{ind\_recibo\_ult1}): 1,22\% train,
  1,07\% validation. $\lambda=\lambda_1=10^{-3}$.}
\end{frame}

\begin{frame}{Dữ liệu và cỡ mẫu}
  \begin{itemize}
    \item \textbf{Đã giảm cỡ mẫu để huấn luyện các thuật toán dễ dàng hơn}: chỉ giữ
          \textbf{1/12 số khách hàng} (lọc theo mã khách, \emph{cùng một} tập khách xuyên suốt
          mọi tháng để ghép lịch sử vẫn khớp), thay vì toàn bộ $\sim$13,6 triệu dòng.
          Nhờ vậy ma trận thiết kế vừa RAM và mỗi thuật toán chạy trong vài giây đến vài chục giây.
    \item Cửa sổ thời gian: 13 tháng, 2015-01 $\to$ 2016-01. Validation là tháng \textbf{2016-01}.
    \item Mỗi mẫu cần đủ \textbf{5 tháng lịch sử}, nên tháng mục tiêu của tập train chỉ từ
          2015-06 đến 2015-12.
    \item \textbf{Cùng một mẫu và cùng bộ đặc trưng} cho mọi mô hình (các thuật toán tự viết,
          sklearn, LightGBM) $\Rightarrow$ so sánh công bằng giữa các mô hình; nhưng
          \emph{không} so trực tiếp được với điểm Kaggle (chấm trên toàn bộ khách, tháng 2016-06).
  \end{itemize}
\end{frame}

\begin{frame}{160 đặc trưng --- theo hướng các lời giải top Kaggle}
  \begin{center}
  \footnotesize
  \begin{tabular}{llp{7.6cm}}
    \toprule
    \textbf{Nhóm} & \textbf{Số cột} & \textbf{Cách tính} \\
    \midrule
    Chặn            & 1  & cột toàn số 1, \textbf{không bị phạt} \\
    Số              & 6  & \texttt{age}, \texttt{antiguedad}, \texttt{renta} ($\log(1{+}x)$),
                           \texttt{ind\_nuevo}, \texttt{ind\_actividad\_cliente}, \texttt{alta\_year} \\
    Nhị phân        & 7  & one-hot bỏ mức tham chiếu (giới tính, phân khúc, quan hệ, 3 kênh phổ biến nhất) \\
    Sở hữu $t-1$    & 24 & đã có sản phẩm đó ở tháng $t-1$ hay chưa \\
    Sở hữu $t-2..t-5$ & 96 & như trên, lùi 2--5 tháng (chưa mở tài khoản $=$ chưa có) \\
    Thời gian       & 24 & số tháng kể từ lần gần nhất có sản phẩm (1--5; 6 $=$ không có trong 5 tháng) \\
    Tổng hợp        & 3  & số sản phẩm đang có; phân khúc / trạng thái hoạt động có đổi so với $t-1$ \\
    \bottomrule
  \end{tabular}
  \end{center}
  \vfill
  Mọi cột chuẩn hoá $(x-\mu)/\sigma$ theo $\mu,\sigma$ \textbf{của tập train} --- nếu không,
  \texttt{renta} cỡ vạn còn cờ 0/1 cỡ 1, gradient descent sẽ zig-zag rất lâu.
\end{frame}

% =====================================================================
\section{GD với bước cố định}

\begin{frame}{GD: chọn độ dài bước thế nào?}
  \[ w_{k+1}=w_k-t\,\nabla F(w_k) \]
  Thử $t \in \{0.1,\,0.5,\,1,\,1.5,\,2,\,4\}\times \tfrac1L$ ($4/L$ vượt ngưỡng $2/L$ của lý thuyết).
  \begin{columns}[T]
    \column{0.5\textwidth}\hinh{gd_codinh_vonglap.pdf}
    \column{0.5\textwidth}\hinh{gd_codinh_thoigian.pdf}
  \end{columns}
\end{frame}

\begin{frame}{Kinh nghiệm rút ra --- bước cố định}
  \begin{itemize}
    \item $t=1/L$ \textbf{luôn hội tụ} nhưng quá bi quan: sau 250 vòng $F=0.0671$, còn xa
          sklearn ($0.0436$).
    \item Bước càng lớn càng tốt: $2/L\to0.0552$, $4/L\to\mathbf{0.0482}$ (tốt nhất).
    \item \textbf{$4/L$ vượt ngưỡng $2/L$ mà vẫn hội tụ}, nhưng \emph{dao động} ở $\sim$10 vòng đầu
          ($F$ tăng lại vài lần) rồi mới đi xuống đều:
      \begin{itemize}
        \item tại $w_0=0$ mọi $\sigma(z_i)=\frac12$ nên $\sigma'=\frac14$ --- chặn $L$ \emph{đúng}
              ở đây, bước $4/L$ vọt lố theo hướng cong nhất;
        \item càng về gần nghiệm, xác suất dự đoán $\approx1\%$, $\sigma'\approx0.01\ll\frac14$:
              độ cong thật nhỏ hơn $L$ nhiều lần, bước $4/L$ lại trở thành an toàn.
      \end{itemize}
    \item[$\Rightarrow$] $2/L$ là ngưỡng cho \emph{trường hợp xấu nhất}; ngưỡng thật phụ thuộc
          vùng mà thuật toán đang đi qua --- lý do backtracking đáng giá.
  \end{itemize}
\end{frame}

% =====================================================================
\section{GD với backtracking}

\begin{frame}{Backtracking (Armijo)}
  Mỗi vòng thử bước \textbf{gấp đôi} lần trước, chưa thoả thì nhân với $\beta$:
  \[ F(w-t\nabla F)\;\le\;F(w)-c\,t\,\|\nabla F\|^2,\qquad c=10^{-4} \]
  \alert{Phải cho bước lớn dần.} Nếu chỉ biết co xuống từ $1/L$ thì backtracking vô dụng ---
  $1/L$ vốn đã luôn thoả điều kiện nên sẽ không bao giờ đổi gì.
  \begin{columns}[T]
    \column{0.5\textwidth}\hinh{gd_bt_vonglap.pdf}
    \column{0.5\textwidth}\hinh{gd_bt_thoigian.pdf}
  \end{columns}
\end{frame}

\begin{frame}{Kinh nghiệm rút ra --- backtracking}
  \begin{itemize}
    \item Theo \textbf{số vòng lặp}: mọi $\beta$ gần như trùng nhau, $F\in[0.04363,\,0.04369]$
          --- $\beta$ không đổi hướng đi.
    \item Theo \textbf{thời gian}: tách hẳn ra. $\beta=0.1$ mất 36 giây, $\beta=0.9$ mất
          140 giây ($\approx3{,}9$ lần) vì mỗi vòng phải tính $F$ nhiều lần mới co đủ.
    \item So với bước cố định tốt nhất ($4/L$, $F=0.0482$): backtracking xuống tới $0.04363$,
          \emph{thấp hơn cả sklearn}, không cần biết $L$.
    \item[$\Rightarrow$] Đây chính là lý do đề bài bắt vẽ \textbf{hai hình riêng}: một cấu hình có
          thể thắng ở hình này và thua ở hình kia.
  \end{itemize}
\end{frame}

% =====================================================================
\section{GD tăng tốc (Nesterov)}

\begin{frame}{Nesterov}
  \[
    w_{k+1}=v_k-t\nabla F(v_k),\qquad
    v_{k+1}=w_{k+1}+\tfrac{\theta_k-1}{\theta_{k+1}}(w_{k+1}-w_k)
  \]
  Cùng thông tin với GD (chỉ cần $L$), tốc độ $O(1/k^2)$ thay vì $O(1/k)$.
  \begin{columns}[T]
    \column{0.5\textwidth}\hinh{nesterov_vonglap.pdf}
    \column{0.5\textwidth}\hinh{nesterov_thoigian.pdf}
  \end{columns}
  \vfill
  {\footnotesize Cùng bước $1/L$, cùng 250 vòng: GD $F=0.0671$, Nesterov $F=0.04368$.
  Tốt nhất ở $t=2/L$: $F=0.04364$ (33 giây).}
\end{frame}

% =====================================================================
\section{Newton}

\begin{frame}{Newton}
  \[
    \nabla^2F(w_k)\,d_k=\nabla F(w_k),\qquad w_{k+1}=w_k-t_k\,d_k,\qquad
    \nabla^2F=\tfrac1nX^\top\mathrm{diag}\big(\sigma_i(1-\sigma_i)\big)X+\lambda I'
  \]
  \begin{columns}[T]
    \column{0.5\textwidth}\hinh{newton_codinh_vonglap.pdf}
    \column{0.5\textwidth}\hinh{newton_bt_vonglap.pdf}
  \end{columns}
  \vfill
  {\footnotesize Trái: bước cố định $t\in\{0.1,0.25,0.5,1\}$. Phải: backtracking (thử $t=1$
  trước), $\beta\in\{0.1,0.5,0.9\}$. 30 vòng mỗi cấu hình.}
\end{frame}

\begin{frame}{Kinh nghiệm rút ra --- Newton}
  \begin{itemize}
    \item $t=0.5$ và $t=1$ đều tới $F=\mathbf{0.043630}$ --- giá trị \textbf{thấp nhất} trong mọi
          thuật toán, trong $\sim$10 giây.
    \item $t=0.1$ (damped quá mạnh) chỉ tới $0.0624$: bỏ mất hội tụ bậc hai.
    \item Backtracking: cả ba $\beta$ cùng tới $0.043630$ (gần nghiệm $t=1$ luôn được nhận),
          nhưng $\beta=0.9$ mất 129 giây vì ở các vòng đầu phải co rất nhiều lần.
    \item Mỗi vòng đắt ($O(nd^2)$ dựng Hessian, $d=161$) nhưng chỉ cần vài vòng
          $\Rightarrow$ vẫn nhanh nhất theo thời gian trong các thuật toán tự viết.
  \end{itemize}
\end{frame}

% =====================================================================
\section{SGD}

\begin{frame}{SGD --- mini-batch}
  \[ \nabla\ell(w)\;\approx\;\frac1B\sum_{i\in\mathcal{B}}\nabla\ell_i(w) \]
  Ước lượng \textbf{không chệch}, phương sai $\sigma^2/B$; rẻ hơn $n/B$ lần mỗi bước.
  \begin{columns}[T]
    \column{0.5\textwidth}\hinh{sgd_eta_epoch.pdf}
    \column{0.5\textwidth}\hinh{sgd_eta_thoigian.pdf}
  \end{columns}
  \vfill
  {\footnotesize Trục tung là \textbf{loss trên mini-batch} (20 lần ghi mỗi epoch) nên đường gợn
  đúng bản chất nhiễu của SGD. Trục hoành đếm theo \textbf{epoch} (1 epoch $=$ 1 lượt đọc hết
  dữ liệu $\approx$ chi phí 1 vòng GD).}
\end{frame}

\begin{frame}{SGD --- lịch teo bước}
  Điều kiện Robbins--Monro:
  $\sum_k\eta_k=\infty$ (còn đủ đà tới đáy), $\sum_k\eta_k^2<\infty$ (nhiễu tắt dần).
  \begin{columns}[T]
    \column{0.5\textwidth}\hinh{sgd_lich_epoch.pdf}
    \column{0.5\textwidth}\hinh{sgd_lich_thoigian.pdf}
  \end{columns}
\end{frame}

\begin{frame}{Loss mini-batch và $F$ trên toàn bộ dữ liệu}
  \begin{columns}[T]
    \column{0.55\textwidth}\hinh{sgd_muot_vs_batch.pdf}
    \column{0.45\textwidth}
      \begin{itemize}
        \item Bước cập nhật luôn nhiễu (mỗi bước chỉ nhìn $B/n\approx1\%$ dữ liệu).
        \item Nhưng $F$ đo trên \textbf{toàn bộ} tập tại $w_k$ là số \textbf{tất định} ---
              vẽ nó thì đường SGD mượt như GD.
        \item Đường gợn hay mượt là do \emph{đại lượng đem vẽ}, không phải do đường đi của $w_k$.
      \end{itemize}
  \end{columns}
\end{frame}

\begin{frame}{Kinh nghiệm rút ra --- SGD}
  \begin{itemize}
    \item Bước cố định, 42 epoch: $\alpha_0=0.2$ tốt nhất (loss epoch cuối $0.04423$);
          $\alpha_0=0.01$ quá nhỏ ($0.0633$), $\alpha_0=5$ quá nhiễu ($0.0474$).
    \item \textbf{Backtracking không dùng được cho SGD}: điều kiện Armijo cần so $F$ trước/sau,
          mà $F$ trên mini-batch là số nhiễu nên sẽ chấp nhận/từ chối bừa.
    \item $B=1$ sẽ \emph{chậm hơn} GD theo giây: một bước GD là \textbf{một} lệnh BLAS trên ma
          trận lớn, còn SGD là hàng trăm lệnh nhỏ. Đo được: 1 epoch ($\approx0{,}96$ s) tốn gấp
          $\sim$15 lần 1 vòng GD ($\approx0{,}064$ s) dù $B=4096$.
    \item Teo bước \textbf{chỉ đáng khi chạy đủ lâu}. Ở 42 epoch, $\alpha_0/k$ teo quá nhanh
          ($0.0527$), $\alpha_0/\sqrt k$ được $0.0459$, còn bước cố định thắng ($0.0442$).
  \end{itemize}
\end{frame}

% =====================================================================
\section{Subgradient và Proximal cho Lasso}

\begin{frame}{Subgradient}
  $|w_j|$ không khả vi tại 0, thay gradient bằng phần tử \textbf{chuẩn nhỏ nhất} của dưới vi phân:
  \[
    \partial|w_j|=\begin{cases}\{\operatorname{sign}(w_j)\} & w_j\neq0\\ [-1,1] & w_j=0\end{cases}
    \qquad
    s_j=\nabla_j\ell-\operatorname{clip}\big(\nabla_j\ell,-\lambda_1,\lambda_1\big)
  \]
  Khi $|\nabla_j\ell|\le\lambda_1$ thì $s_j=0$: toạ độ \textbf{đứng yên tại 0} $\Rightarrow$ nghiệm thưa.
  \begin{columns}[T]
    \column{0.5\textwidth}\hinh{subgrad_vonglap.pdf}
    \column{0.5\textwidth}\hinh{subgrad_thoigian.pdf}
  \end{columns}
\end{frame}

\begin{frame}{Proximal gradient (ISTA)}
  \[
    w_{k+1}=\operatorname{prox}_{t\lambda_1\|\cdot\|_1}\big(w_k-t\nabla\ell(w_k)\big),\qquad
    \operatorname{prox}(v)_j=\operatorname{sign}(v_j)\max(|v_j|-t\lambda_1,0)
  \]
  Tách phần trơn và phần $L_1$: bước cố định $t\le1/L$ dùng được, soft-threshold đặt \textbf{hẳn}
  hệ số về 0.
  \begin{columns}[T]
    \column{0.5\textwidth}\hinh{proximal_vonglap.pdf}
    \column{0.5\textwidth}\hinh{proximal_thoigian.pdf}
  \end{columns}
\end{frame}

\begin{frame}{Kinh nghiệm rút ra --- Lasso}
  \begin{itemize}
    \item \textbf{Subgradient không phải phương pháp giảm}: $F$ vọt lên ở vài vòng đầu rồi mới
          xuống --- đường gợn, khác hẳn GD. Bước bắt buộc teo dần, tốc độ chỉ $O(1/\sqrt{k})$.
          Tốt nhất $\alpha_0=8$: $F=0.0517$, 154/160 hệ số khác 0 (gần như không thưa).
    \item \textbf{Proximal thưa thật sự}: $t=2/L$ cho 30/160 hệ số khác 0.
          Nhưng với $d=161$ và 250 vòng, ISTA \emph{chưa hội tụ}: $F=0.0583$, còn cao hơn
          subgradient --- tốc độ $O(1/k)$ với $L=6.25$ là chậm; cần FISTA hoặc thêm vòng.
    \item sklearn \texttt{liblinear}: $F=0.0467$, 19/160 hệ số khác 0 --- cả hai code tự viết đều
          chưa tới nghiệm Lasso trong 250 vòng.
  \end{itemize}
\end{frame}

% =====================================================================
\section{So sánh và kết luận}

\begin{frame}{Năm thuật toán (Ridge) ở setup tốt nhất, so với sklearn}
  \begin{columns}[T]
    \column{0.5\textwidth}\hinh{tongket_vonglap.pdf}
    \column{0.5\textwidth}\hinh{tongket_thoigian.pdf}
  \end{columns}
  \vfill
  {\footnotesize Điều kiện để so được: cùng tối thiểu hoá \emph{một} hàm $F$ (Ridge, cùng
  $\lambda$). Quy đổi cho sklearn: nó dùng $\frac12\|w\|^2+C\sum_i\ell_i$ (tổng), ta dùng trung
  bình $\Rightarrow C=\dfrac{1}{n\lambda}$. Trục hoành của SGD tính theo epoch.}
\end{frame}

\begin{frame}{Lasso: subgradient và proximal so với thư viện}
  \begin{columns}[T]
    \column{0.5\textwidth}\hinh{lasso_vonglap.pdf}
    \column{0.5\textwidth}\hinh{lasso_thoigian.pdf}
  \end{columns}
  \vfill
  {\footnotesize Để riêng vì Lasso tối thiểu hoá \emph{hàm khác} ($\lambda_1\|w\|_1$ thay
  cho $\frac\lambda2\|w\|_2^2$) --- đặt chung một trục với Ridge là so sai.}
\end{frame}

\begin{frame}{Bảng kết quả}
  \begin{center}
  \footnotesize
  \begin{tabular}{lrr}
    \toprule
    \textbf{Setup tốt nhất (Ridge)} & \textbf{$F$ cuối} & \textbf{Thời gian (giây)} \\
    \midrule
    GD bước cố định $t=4/L$              & 0.048206 & 15.2 \\
    GD backtracking $\beta=0.1$          & 0.043634 & 36.0 \\
    GD Nesterov $t=2/L$                  & 0.043636 & 33.3 \\
    Newton $t=1$ (30 vòng)               & \textbf{0.043630} & 10.2 \\
    SGD $\alpha_0=0.2$, bước cố định$^*$ & 0.044234 & 40.4 \\
    \midrule
    sklearn \texttt{LogisticRegression} (mặc định, lbfgs) & 0.043639 & 2.0 \\
    \bottomrule
  \end{tabular}

  \vspace{0.5em}
  \begin{tabular}{lrrr}
    \toprule
    \textbf{Lasso} & \textbf{$F$ cuối} & \textbf{Giây} & \textbf{Hệ số $\neq 0$} \\
    \midrule
    Subgradient $\alpha_0=8$                           & 0.051692 & 14.4 & 154/160 \\
    Proximal (ISTA) $t=2/L$                            & 0.058274 & 14.6 & 30/160 \\
    sklearn \texttt{liblinear} (\texttt{penalty='l1'}) & 0.046689 & 9.7  & 19/160 \\
    \bottomrule
  \end{tabular}
  \end{center}
  {\scriptsize $^*$ SGD: trung bình loss mini-batch ở epoch cuối, không phải $F$ trên toàn bộ.}
\end{frame}

% =====================================================================
\section{Mở rộng: 24 sản phẩm và so với top Kaggle}

\begin{frame}{MAP@7 --- logistic so với mô hình của các lời giải top Kaggle}
  24 sản phẩm, cùng mẫu, cùng 160 đặc trưng, validation 2016-01. Chỉ đổi \textbf{mô hình}.
  \begin{center}
  \footnotesize
  \begin{tabular}{lrrr}
    \toprule
    \textbf{Mô hình} & \textbf{MAP@7 khách có thêm} & \textbf{MAP@7 toàn bộ} & \textbf{Giây} \\
    \midrule
    Phổ biến (mốc sàn)                                        & 0.6264 & 0.01768 & 0 \\
    Logistic L1 $\times$24 (Proximal, $t=2/L$)                & 0.7733 & 0.02183 & 258 \\
    LightGBM nhị phân $\times$24 {\scriptsize(hướng hạng 2)}  & 0.8087 & 0.02283 & 44 \\
    LightGBM softmax {\scriptsize(hướng hạng 1)}              & \textbf{0.8197} & \textbf{0.02314} & 8 \\
    \bottomrule
  \end{tabular}
  \end{center}
  \begin{itemize}
    \item Đổi sang 160 đặc trưng: logistic tăng từ $0.7467\to0.7733$ (bộ 38 cột cũ).
    \item Cùng đặc trưng, LightGBM softmax hơn logistic $\sim$4,6 điểm --- nhiều khả năng vì cây
          bắt được tương tác giữa các đặc trưng (vd.\ "có sản phẩm A tháng trước \emph{và} mới đổi
          phân khúc") mà mô hình tuyến tính không biểu diễn được.
    \item Kaggle hạng 1 $\approx0.0314$ (toàn bộ khách, tháng 2016-06, đặc trưng lớn hơn) --- không so
          trực tiếp được.
  \end{itemize}
\end{frame}

\begin{frame}{Kết luận}
  \begin{enumerate}
    \item \textbf{Về giá trị hàm mục tiêu}: Newton ($0.043630$), GD backtracking ($0.043634$) và
          Nesterov ($0.043636$) đều \textbf{thấp hơn} sklearn lbfgs mặc định ($0.043639$) ---
          lbfgs dừng sớm khi đạt ngưỡng sai số mặc định (\texttt{tol}$=10^{-4}$).
    \item \textbf{Về thời gian}: sklearn vẫn nhanh nhất (2 giây); Newton tự viết chậm hơn
          $\sim$5 lần, GD backtracking $\sim$18 lần.
    \item \textbf{Chọn tham số}: bước cố định phải dò tay và $2/L$ chỉ là ngưỡng xấu nhất ($4/L$
          vẫn hội tụ); backtracking tự tìm bước mà không cần biết $L$.
    \item \textbf{Số vòng lặp $\neq$ thời gian}: mọi kết luận đều phải đọc trên cả hai hình.
    \item \textbf{Điều kiện của bài toán}: $L=6.25$, $\mu\ge\lambda=10^{-3}$ nên $L/\mu$ có thể tới
          $\sim$6000 --- GD cần rất nhiều vòng; đó là lý do các phương pháp gia tốc
          ($O(\sqrt{L/\mu})$) và Newton (bậc hai) thắng rõ.
    \item \textbf{Về mô hình}: với cùng đặc trưng, gradient boosting vượt logistic trên MAP@7 ---
          tối ưu tốt hàm log-loss tuyến tính không bù được giới hạn của mô hình tuyến tính.
  \end{enumerate}
\end{frame}

\end{document}
